\section{Résultats}

\subsection{Oscillations libres}

Avant de mesurer une oscillation, nous avons mesure le rapport de force qu'il faut exercer sur le balancier afin de le placer a un certain angle.
\\
\noindent \begin{minipage}{0.5\textwidth}
\includegraphics[width=\textwidth]{raideur.png}
	\captionof{figure}{$M_r$($\theta$)}
\label{fig:C}
\end{minipage}
\begin{minipage}{0.4\textwidth}
	Nous obtenons par consequent le graphique (\ref{fig:C}) ci-joint.
	Grace a sa pente d'unite ... nous avons pu directement extraire la valeur C, soit la constante de raideur du ressort utilise grace a la relation (\ref{eq:nature}).\\ Soit :
	\\      C = 0.0123 +/- 0.0002 [Nm / $rad^-1$]
\noindent \end{minipage}
\\
\\
\\
En un second temps, apres avoir montee l'installation pour mesures les oscillations du balencier-spiral (voir fig:\ref{fig:samplesetup}), nous avons donne une position hors-equilibre au balencier et mesurer en fonction du frein imposee a les differences oscillations observable en annexe (voir fig:\ref{fig:oscillations}).

\begin{minipage}[t]{0.45\textwidth}
    \centering
    \includegraphics[width=\linewidth]{lambda_quadratique.png}
	\captionof{figure}{$\lambda$(I)}
	\label{fig:lambda-current_quad}
\end{minipage}
\hfill
\begin{minipage}[t]{0.5\textwidth}
    \centering
    \includegraphics[width=\linewidth]{lambda.png}
	\captionof{figure}{$\lambda$($I^2$)}
	\label{fig:lambda-current}
\end{minipage}
\\
\\
La relation entre le facteur d'amortissement et le courant circulant dans les freins n'est pas lineaire. Nous constatons verifions la relation (\ref{eq:lambda-courant}), avec le graphique \ref{fig:lambda-current}. Nous pouvons extraire grace a la relation \ref{eq:JCW} la moyenne des pulsations propres $\overline{\omega_0}$ = 3.25 +/- 0.02 [rad/s] (voir tab:\ref{val:omega0}) desquelles on extrait la frequence propre : $f_0$ = ... (a partir de l'inverse de la periode : Eq. \ref{eq:t}). Avec la constante de torsion du ressort le moment d'inertie du balencier $J$ = $C$ / $\omega_^2$ = $0.0123/3.25^2$ = (1.16 +/- 0.000) \cdot $10^3$  [kg / $m^2$]. 

\newpage
\subsection{Expérience 2}

